\documentclass[11pt]{article}
\textwidth 15.9cm
\textheight 23. cm
\addtolength{\oddsidemargin}{-11.5mm}
\addtolength{\evensidemargin}{-11.5mm}
\addtolength{\topmargin}{-19.0mm}

%\usepackage{amsmath}
\usepackage[T1]{fontenc}
\usepackage[ansinew]{inputenc}
\usepackage[bbgreekl]{mathbbol}
\usepackage{amsmath,amssymb}
\usepackage{graphicx}
\usepackage[colorlinks=true,urlcolor=cyan]{hyperref}%,linkcolor=green]{hyperref}
\usepackage{url}
\usepackage[numbers,square]{natbib}
\usepackage{multirow}
\usepackage{tikz}
\usepackage{tikz-cd}
\usetikzlibrary{arrows,matrix}
\usepackage{bbm}
\usepackage{authblk}
\usepackage{subcaption}
\usepackage{xstring}
\usepackage{xparse}
\usepackage{comment}
\renewcommand{\Affilfont}{\footnotesize}
\usepackage{stmaryrd}
\usepackage{comment}
\usepackage{showlabels}
\usepackage{soul}

\usetikzlibrary{calc,patterns,shapes.geometric}
\usetikzlibrary{arrows,shapes,positioning}
\usetikzlibrary{decorations.markings,decorations.pathmorphing,decorations.pathreplacing}
\usetikzlibrary{decorations.text,decorations.shapes}

%%%%%%%% MACROS
\input{GempicLatexMacros}

\title{Diagnostics}


\author[1]{The Gempic development team}
\affil[1]{Max-Planck-Institut f\"ur Plasmaphysik, Garching, Germany}


\begin{document}

\maketitle

\begin{abstract}
This note describes the use of the two Diagnostics modules, FullDiagnostics and ReducedDiagnostics.
\end{abstract}

\tableofcontents
%
GEMPIC has two modules for outputting diagnostics: FullDiagnostics for outputting entire fields and ReducedDiagnostics for outputting metrics based on the fields.

In simulation code a \lstinline|Diagnostics| object must be initialised in the main code, and a call to
\lstinline|Diagnostics::compute_and_write_to_file|\\
must then be made for each time step. This will automatically call both Diagnostics modules with the correct parameters.

\section{The FullDiagnostics Module}
FullDiagnostics produces one or more Paraview-readable AMReX-format output files with the binary data of the fields asked for in the input file (see the corresponding section in Input Files).
Assuming the module is enabled in the input file, full diagnostic groups can be named and base parameters can be specified. Each diagnostics group can either treat particles or fields.

Where particle positions are simply printed in binary, for fields there are a number of output processors to choose from, all selected via the line(s)
\\\lstinline|FullDiagnostics.{groupName}.outputProcessor|\\
in the input file:

\begin{itemize}
    \item \texttt{"Raw"}: The "\emph{raw}" field output, which is simply the field copied over with no processing. Due to the way Paraview works, the field is assumed to be node centered.
    \item \texttt{"CellCenter"}: The interpolated field output for cell centered fields. Does a very simple linear interpolation to every node from all the nearest cell centers surrounding it.
\end{itemize}
If more advanced functionality is needed the user can also either write their own custom operator (\ref{sec:cusOp}) or custom class (\ref{sec:cusClass}), selecting the \verb|"Custom"| output processor and the appropriate \verb|customID|.
Both these options require some extra work.

\subsection{Custom Operators}\label{sec:cusOp}
Users can add their own operators \emph{before the first instantiation of} \lstinline|FullDiagnostics| to be called on the GPU to calculate the output fields from the input fields specified by name in the input file. As such, they must have a very specific signature and be explicitly added to the list of custom processor options in the code:
\begin{lstlisting}
std::string id{"NameForCustomOperator"}; // note that ID must be CamelCase
Gempic::Io::add_output_processor(
    id,
    [=] AMREX_GPU_DEVICE(
        amrex::Array4<amrex::Real> dst,             // outputMultifab
        const amrex::Array4<const amrex::Real> src, // inputMultifab
        int nSrcComp,                               // # input multifab components
        int i, int j, int k,                        // indices
        double scaling,                             // scaling and shifts
        double ishift, double jshift, double kshift)// for cell centered fields
    {
        ...
    });
\end{lstlisting}
Note that even if your custom operator is only meant to work for node centered fields, it still must accept the values \verb|scaling|, \verb|ishift|, \verb|jshift| and \verb|kshift| needed for cell-centered fields.
Also note that the lambdas of course can capture copyable objects to use for their functionality.

Then, the custom operators are selected in the input file using the id given in the code:
\\\lstinline|FullDiagnostics.{groupName}.customID = "NameForCustomOperator"|\\
%
For an example, see \verb|Testing/Io/FullDiagnostics_test.cpp|.

\subsection{Custom Classes}\label{sec:cusClass}
If an operator is not enough, users can also add their own custom class. This class must
\begin{itemize}
    \item inherit from \verb|Gempic::Io::CustomOutputProcessor|,
    \item implement a public operator with the signature\\
    \lstinline|void operator()(amrex::MultiFab& mfDst, int dcomp) const override|\\or\\
    \lstinline|void operator()(amrex::MultiFab& mfDst, int dcomp) const final|,
    where \lstinline|mfDst| is the output MultiFab to be written and \lstinline|dcomp| is the first component of \lstinline|mfDst| in which data is stored.
\end{itemize}

Finally, the custom class must be added in the code:
\begin{lstlisting}
    class YourOutputProcessor : public CustomOutputProcessor {...};
    std::string id{"MyCustomClass"}; // note that ID must be CamelCase
    Gempic::Io::add_output_processor<YourOutputProcessor>(id);
\end{lstlisting}
and selected in the input file:
\\\lstinline|FullDiagnostics.{groupName}.customID = "MyCustomClass"|\\
For an example, see \verb|Testing/Io/FullDiagnostics_test.cpp|.

\section{The ReducedDiagnostics Module}
ReducedDiagnostics produces one or more plaintext output files with the metrics asked for in the input file (see the corresponding section in Input Files).
Assuming the module is enabled in the input file, reduced diagnostic groups can be named and base parameters can be specified. Each diagnostics group can treat both particles and fields.

There are a number of outputs to choose from, all selected via the line(s) 
\\\lstinline|ReducedDiagnostics.{groupName}.types|\\
in the input file:

\begin{itemize}
    \item \texttt{Particle}: The particle metrics, which for now comprises the three dimensional momentum and total kinetic energy of all particles for each particle type.
    \item \texttt{ElecFieldEnergy}: The electric field energy $\langle \arr{E}, \tilde{\arr{D}}\rangle/2$
    \item \texttt{MagFieldEnergy}: The magnetic field energy $\langle \arr{B}, \tilde{\arr{H}}\rangle/2$
    \item \texttt{GaussError}: The Gauss error $||\tilde{\arr{\rho}}-\Div \tilde{\arr{D}}||/||\tilde{\arr{\rho}}||$
    \item \texttt{CurrentFieldEnergy}: The energy of a current field $\langle \arr{J}, \tilde{\arr{J}}\rangle/2$ where the primal and dual fields are related by the density $\tilde{\arr{J}} = n(x) \arr{J}$. If \lstinline|computeMissingFields| is used, the density has to be provided as function \lstinline|FieldDensity|/\lstinline|FieldDensityInv| in the input file.
\end{itemize}
It is currently not possible to provide custom expansions to this list.
Note that multiple types can be selected for each reduced diagnostics group.
There are two general modes for the ReducedDiagnostics module:
\begin{itemize}
    \item The user is responsible for keeping every field needed for their ReducedDiagnostics calculations up to date whenever \lstinline|ReducedDiagnostics::write_to_file| is called. This is the default.
    \item The user does not need to provide all fields required for the ReducedDiagnostics calculations. The ReducedDiagnostics module will attempt to calculate the missing fields whenever \lstinline|Diagnostics::compute_and_write_to_file| is called. This may impact computation time, but can be useful when the fields in question are not needed for the main simulation.
\end{itemize}
The modes can be toggled via the input file switch
\\\lstinline|ReducedDiagnostics.computeMissingFields|\\
in general, and/or
\\\lstinline|ReducedDiagnostics.{groupName}.computeMissingFields|\\
for each group.
\end{document}
